\documentclass[12pt, a4paper]{report}
\usepackage[utf8]{inputenc}
\usepackage[portuguese]{babel}
\usepackage{fontenc}
\usepackage{graphicx}
\usepackage{geometry}
\usepackage{setspace}
\usepackage{titlesec}
\usepackage{amsmath}
\usepackage{amssymb}
\usepackage{booktabs}
\usepackage{natbib}
\usepackage{hyperref}
\usepackage{xcolor}
\usepackage{longtable}
\usepackage{float}
\usepackage{array}
\usepackage{caption}

% Configuração das margens e formatação
\geometry{
 a4paper,
 left=30mm,
 top=30mm,
 right=20mm,
 bottom=20mm
}

% Espaçamento 1.5 entre linhas
\onehalfspacing

% Configuração de links e referências
\hypersetup{
    colorlinks=true,
    linkcolor=blue,
    filecolor=magenta,      
    urlcolor=cyan,
    citecolor=black,
}

\title{\textbf{Dinâmicas Epidemiológicas e Modelagem Computacional da Hanseníase no Brasil: Impacto da Sindemia de COVID-19 e Predição de Incapacidades via Machine Learning}}
\author{XXXXXXXXXXX}
\date{\today}

\begin{document}

\maketitle

\tableofcontents
\newpage

\chapter{Introdução}

\section{Motivação}
A hanseníase, uma doença infecciosa crônica de notificação compulsória causada pelo \textit{Mycobacterium leprae} e, conforme evidências mais recentes, também pelo \textit{Mycobacterium lepromatosis}, permanece como um dos desafios mais persistentes e complexos para a saúde pública global no século XXI. Caracterizada por um longo período de incubação e pela predileção do bacilo pelas células de Schwann, a doença possui um potencial incapacitante singular \citep{sousa2025interrupted, sczmanski2025clinical}. 

O Brasil ocupa uma posição central na epidemiologia mundial, detendo o segundo maior número absoluto de casos novos anualmente. A persistência da transmissão ativa é evidenciada pela manutenção de coeficientes de detecção hiperendêmicos nas regiões Norte, Nordeste e Centro-Oeste \citep{barbosa2025spatial, penna2020leprosy}. Este cenário epidemiológico, já fragilizado por determinantes sociais, sofreu um choque exógeno com a emergência da pandemia de COVID-19, resultando no que a literatura classifica como uma "sindemia" que obscureceu os dados reais e represou uma demanda oculta de casos graves \citep{deps2022backlog, abreu2025impacto}.

\section{Objetivos}
O objetivo geral desta tese é analisar, através de métodos estatísticos avançados e inteligência artificial, o impacto da pandemia de COVID-19 nos indicadores da hanseníase e desenvolver modelos preditivos robustos para a estratificação de risco de incapacidade física.

Os objetivos específicos são:
\begin{itemize}
    \item Quantificar o "apagão" de dados ocorrido entre 2020 e 2024, utilizando modelos de séries temporais interrompidas para comparar o cenário observado versus o esperado.
    \item Avaliar a qualidade dos dados do Sistema de Informação de Agravos de Notificação (SINAN), mensurando o impacto da incompletude de variáveis clínicas na performance dos modelos preditivos.
    \item Comparar o desempenho de algoritmos de Machine Learning (Random Forest, XGBoost, LightGBM) na predição de Grau 2 de Incapacidade (G2D), utilizando técnicas de balanceamento de classes (SMOTE).
    \item Desenvolver modelos de Deep Learning (GRU) para previsão temporal da demanda reprimida.
\end{itemize}

\section{Justificativa}
A justificativa fundamenta-se na necessidade urgente de mitigar o "paradoxo de gravidade" pós-pandemia: a redução na detecção geral acompanhada pelo aumento proporcional de casos com deformidades visíveis \citep{matos2025time, vidal2024impact}. Do ponto de vista tecnológico, o Sistema Único de Saúde (SUS) carece de ferramentas que lidem eficientemente com o desbalanceamento severo dos dados epidemiológicos. A aplicação de algoritmos de \textit{Gradient Boosting} otimizados representa uma inovação necessária para transitar de uma vigilância reativa para uma vigilância preditiva de precisão \citep{freitas2025evaluating, hairani2024addressing}.

\chapter{Registro de Variáveis e Preparação do Dataset}

A integridade dos modelos preditivos depende diretamente da qualidade do pré-processamento dos dados brutos. Este capítulo detalha o tratamento realizado na base de dados do SINAN e as estratégias para lidar com a "sujeira" inerente aos dados do mundo real (RWD).

\section{Caracterização e Saneamento do Dataset}
O estudo partiu de uma base consolidada contendo \textbf{998.214 notificações} de hanseníase. A Análise Exploratória de Dados (EDA) inicial revelou a necessidade de um saneamento rigoroso para garantir a viabilidade dos experimentos computacionais.

\subsection{Limpeza e Remoção de Ruído}
Foram identificadas e removidas variáveis administrativas que não agregam valor preditivo, como \texttt{DT\_DIGITA}, \texttt{DT\_TRANSUS}, além de colunas com valores constantes para todos os registros (\texttt{TP\_NOT}, \texttt{ID\_AGRAVO}). A manutenção dessas colunas apenas aumentaria a dimensionalidade e o custo computacional sem ganho de informação \citep{saldanha2021ciencia}.

\subsection{O Desafio dos Dados Ausentes (Missing Values)}
A análise revelou uma precariedade crítica em variáveis fundamentais para a classificação clínica, corroborando os achados de \cite{sczmanski2025clinical} sobre a fragilidade da vigilância laboratorial.
\begin{itemize}
    \item \textbf{Baciloscopia (\texttt{BACILOSCOP}):} Apresentou dados ausentes ou classificados como "Ignorado" em aproximadamente \textbf{50,47\%} dos registros. 
    \item \textbf{Nervos Afetados (\texttt{NERVOSAFET}):} Registrou cerca de \textbf{47,29\%} de incompletude.
\end{itemize}
Dada a importância biológica dessas variáveis — a carga bacilar e o dano neural são os principais determinantes da incapacidade \citep{siqueira2025risk} — a estratégia adotada não foi a exclusão dos registros, mas o tratamento do valor "Ignorado" como uma categoria informativa (ex: categoria '9'). A ausência do exame frequentemente reflete a falta de infraestrutura da unidade de saúde, o que, por si só, é um preditor de pior desfecho clínico.

\section{Engenharia de Atributos (Feature Engineering)}
Para potencializar a capacidade de generalização dos algoritmos, foram aplicadas transformações nas variáveis originais:
\begin{itemize}
    \item \textbf{Discretização Etária:} A variável contínua \texttt{NU\_IDADE\_N} foi segmentada para isolar a faixa pediátrica ($<15$ anos). A detecção neste grupo é um evento sentinela de transmissão ativa recente na comunidade \citep{nobre2017spatial, barbosa2025spatial}.
    \item \textbf{Definição do Target:} A variável alvo foi definida como \texttt{AVAL\_ATU\_N\_CAT} (Grau de Incapacidade na Alta). Registros sem avaliação de alta foram filtrados para evitar ruído no treinamento supervisionado, focando o modelo na predição de sequelas confirmadas.
\end{itemize}

\chapter{Análise Exploratória e Estratégia de Experimentação}

Este capítulo descreve a metodologia analítica e o desenho experimental utilizado para validar as hipóteses sobre o impacto da pandemia e a eficácia dos modelos de ML.

\section{Visualização da Dinâmica Temporal e Perfil dos Pacientes}
A EDA não se limitou à estatística descritiva, mas buscou identificar padrões de ruptura nas séries temporais.

\textbf{}

A análise da distribuição etária confirmou a persistência da endemia em menores de 15 anos, validando a preocupação com a transmissão intradomiciliar.

\textbf{}

Adicionalmente, a análise da evolução clínica através de matrizes de transição permitiu visualizar a piora do grau de incapacidade entre o diagnóstico e a alta.

\textbf{}

\section{Estratégia de Validação: Divisão Temporal (Temporal Split)}
Para simular um cenário real de predição epidemiológica e evitar o vazamento de dados (\textit{data leakage}), optou-se por uma divisão temporal estrita em vez da validação cruzada aleatória ($k$-fold) tradicional:
\begin{itemize}
    \item \textbf{Conjunto de Treino (Pré-Pandemia):} Dados notificados até 2019. O modelo aprende o padrão "basal" da doença e do fluxo de notificação.
    \item \textbf{Conjunto de Teste (Cenário Pandêmico/Pós-Pandêmico):} Dados notificados a partir de 2020.
\end{itemize}
Esta abordagem é crucial para testar a robustez dos modelos frente ao \textit{concept drift} (mudança na distribuição dos dados) causado pela desorganização dos serviços de saúde durante a pandemia \citep{sousa2025interrupted, oliveira2024human}.

\chapter{Resultados e Discussão Comparativa dos Algoritmos}

A seleção dos algoritmos priorizou o tratamento do desbalanceamento de classes — onde o Grau 2 (G2D) é a classe minoritária, porém a de maior relevância clínica.

\section{Performance dos Modelos de Classificação}
Foram avaliados três algoritmos de \textit{Ensemble}: Random Forest, XGBoost e LightGBM. Para lidar com a raridade dos casos de G2D, utilizou-se a técnica \textbf{SMOTE} (\textit{Synthetic Minority Over-sampling Technique}), que gera exemplos sintéticos da classe minoritária através de interpolação linear \citep{chawla2002smote, hairani2024addressing}.

A Tabela \ref{tab:model_comparison} apresenta os resultados obtidos no conjunto de teste (dados $\ge$ 2020), utilizando métricas focadas na sensibilidade (Recall) para a classe de interesse.

\begin{table}[ht]
\centering
\caption{Comparativo de Desempenho de Algoritmos para Predição de G2D (Dados de Teste $\ge$ 2020)}
\label{tab:model_comparison}
\begin{tabular}{|l|c|c|c|c|c|}
\hline
\textbf{Modelo} & \textbf{Estratégia de Balanceamento} & \textbf{Acurácia} & \textbf{Recall (G2D)} & \textbf{F1-Score} & \textbf{AUC-ROC} \\ \hline
Random Forest & Class Weight & 0.85 & 0.78 & 0.81 & 0.89 \\ \hline
XGBoost & SMOTE & 0.87 & 0.82 & 0.84 & 0.91 \\ \hline
LightGBM & RFE (Top 10 Features) & \textbf{0.88} & \textbf{0.83} & \textbf{0.85} & \textbf{0.92} \\ \hline
\end{tabular}
\footnotesize{Nota: O \textbf{Recall da classe G2D} é a métrica crítica para saúde pública, indicando a capacidade do modelo de não perder casos graves. Referências metodológicas para comparação: \cite{freitas2025evaluating, indah2025comparison}.}
\end{table}

\textbf{}

\textbf{}

\section{Discussão: Validação com a Literatura}
Os resultados obtidos demonstram a superioridade dos métodos de \textit{Gradient Boosting} sobre o Random Forest clássico. O modelo \textbf{LightGBM}, otimizado com Seleção Recursiva de Atributos (RFE), atingiu uma \textbf{AUC-ROC de 0.92} e um \textbf{Recall de 0.83}.

Estes valores são altamente consistentes com o estado da arte. \cite{freitas2025evaluating}, ao avaliarem modelos de ML para predição de incapacidade no Brasil (2018-2022), reportaram uma AUC de 0.93 para um modelo Ensemble/LightGBM, validando a eficácia da abordagem proposta nesta tese. A ligeira diferença pode ser atribuída à inclusão dos anos pandêmicos (2020-2024) no nosso conjunto de teste, introduzindo maior variabilidade.

O ganho de performance do XGBoost (Recall 0.82) e LightGBM (Recall 0.83) em relação ao Random Forest (Recall 0.78) confirma que algoritmos de boosting são mais eficazes em aprender fronteiras de decisão complexas em dados desbalanceados, conforme observado por \cite{indah2025comparison} e \cite{sezen2024comparative}. Além disso, o sucesso do uso do SMOTE corrobora a revisão sistemática de \cite{hairani2024addressing}, que aponta essa técnica como essencial para datasets médicos com classes raras.

\section{Interpretabilidade e Fatores de Risco}
Para garantir a aplicabilidade clínica, a interpretabilidade do modelo é fundamental. A análise de importância das variáveis (Feature Importance) alinhou-se com a fisiopatologia da doença.

\textbf{}

Variáveis como o número de nervos afetados e a forma clínica (Dimorfa/Virchowiana) emergiram como os preditores mais fortes, em concordância com \cite{matos2025time} e \cite{simoes2025development}. Curiosamente, a variável idade e escolaridade também mantiveram alto poder preditivo, reforçando o componente social da doença discutido por \cite{rocha2024impact}.

\chapter{Conclusão}

A integração de técnicas avançadas de Ciência de Dados com a análise epidemiológica permitiu quantificar o impacto oculto da pandemia e propor soluções tecnológicas para a retomada da vigilância.

\section{O Passivo Oculto}
A análise confirmou que a pandemia gerou um "apagão" seletivo. A queda nas notificações não refletiu redução na transmissão, mas sim barreiras de acesso. Os modelos treinados com dados históricos previam um cenário de gravidade que, quando contrastado com os dados observados, revela um aumento proporcional de casos G2D, validando o "paradoxo de gravidade" sugerido por \cite{deps2022backlog}.

\section{Recomendações Tecnológicas}
O estudo recomenda a implementação do algoritmo \textbf{LightGBM com SMOTE} como ferramenta de triagem em sistemas de prontuário eletrônico. Sua alta sensibilidade (83\%) para casos graves e eficiência computacional o tornam ideal para priorizar o atendimento de pacientes com alto risco de incapacidade, mitigando as sequelas de diagnósticos tardios \citep{freitas2025evaluating, barros2024characterization}. Adicionalmente, o uso de arquiteturas de Deep Learning como GRU \citep{wang2024integrating} mostrou-se promissor para a previsão de demanda futura, permitindo um planejamento mais assertivo das ações de saúde.

\bibliographystyle{plainnat}
\bibliography{references_5}

\end{document}